% ======================================================================
% CAPÍTULO 3 — TRABALHOS RELACIONADOS
% ======================================================================
\chapter{TRABALHOS RELACIONADOS}

Este capítulo situa o trabalho desenvolvido em relação a outras abordagens e
sistemas descritos na literatura e na prática de instrumentação analítica,
organizando a comparação em quatro eixos: instrumentação para especiação de
mercúrio, controle de temperatura em fornos analíticos, arquiteturas de
aquisição de dados e plataformas de supervisão.

\section{Instrumentação para especiação de mercúrio}

A literatura de especiação de mercúrio concentra-se sobretudo nos métodos
químicos e nas estratégias de detecção, e menos na engenharia do controle que
automatiza a preparação de amostras. Os trabalhos clássicos de Bloom
\cite{bloom1989} e Liang e colaboradores \cite{liang1994} estabeleceram os
protocolos de etilação em fase aquosa com NaBEt$_4$ e de separação por
cromatografia gasosa com detecção por fluorescência atômica em vapor frio, que
são a base química do processo aqui automatizado. Revisões como a de Leermakers
e colaboradores \cite{leermakers2005} discutem as principais fontes de artefato
e a necessidade de validação dos métodos de especiação, evidenciando que a
repetibilidade das condições de preparação --- em particular da rampa térmica ---
é um fator crítico de qualidade.

No plano instrumental, equipamentos comerciais dedicados, como o analisador de
mercúrio da família RA-915 da Lumex \cite{lumex}, integram a atomização e a
detecção, mas tipicamente não expõem, ao usuário, o controle fino do
sequenciamento da preparação que antecede a medida. Essa observação reforça a
relevância de um sistema de automação dedicado à etapa de preparação, como o
desenvolvido neste trabalho, que pode ser acoplado a detectores comerciais.

\section{Controle de temperatura e rampas em fornos analíticos}

O controle de temperatura em fornos de laboratório é um problema clássico de
controle de processos, frequentemente resolvido com controladores PID
\cite{astrom2006, ogata2010}. A literatura de controle aborda extensamente a
sintonia de controladores PID, o tratamento de saturação (anti-windup) e o
seguimento de setpoints variantes no tempo. No caso específico de rampas de
temperatura, a abordagem de setpoint dinâmico em malha fechada é a mais comum
quando a medição é confiável em toda a faixa. O sistema aqui proposto combina
essa abordagem com uma estratégia de razão de taxas em malha aberta para a
região em que o termopar não fornece leitura confiável --- uma solução de
engenharia motivada pela limitação física do sensor, cuja discussão detalhada é
apresentada no Capítulo~5.

\section{Arquiteturas de aquisição de dados}

No que se refere à arquitetura de hardware, duas famílias de solução são
frequentes na instrumentação analítica: o uso de controladores lógicos
programáveis (CLPs) e de sistemas de aquisição dedicados, tipicamente de custo
elevado, e o uso de plataformas abertas de baixo custo baseadas em
microcontroladores, como o Arduino, e em computadores de placa única, como o
Raspberry Pi \cite{arduino, raspberrypi}. Plataformas abertas têm sido cada vez
mais adotadas em instrumentação científica por sua flexibilidade, custo e
ecossistema de desenvolvimento, especialmente quando os requisitos de tempo real
são moderados e podem ser isolados em um microcontrolador dedicado.

A arquitetura mestre-escravo adotada --- em que um computador de alto nível
executa a lógica de supervisão e um microcontrolador executa a aquisição e o
acionamento de campo em tempo real --- é uma solução de compromisso entre custo
e determinismo. Essa separação também é defendida na engenharia de software
como forma de reduzir o acoplamento entre a lógica de aplicação e os detalhes de
hardware, melhorando a manutenibilidade do sistema.

\section{Plataformas de supervisão e IHM}

Historicamente, sistemas de supervisão laboratoriais eram construídos sobre
plataformas proprietárias e monolíticas, como o LabVIEW, que acoplam a lógica de
controle, a interface e a aquisição em um único ambiente. Embora produtivas para
prototipagem, tais plataformas apresentam limitações de custo de licenciamento,
dificuldade de manutenção e evolução e dependência de um fornecedor. A
modernização para arquiteturas Web, com backend em linguagens de uso geral e
IHM acessível por navegador, permite monitorar o processo a partir de qualquer
dispositivo da rede local, integrar o registro de dados a sistemas de auditoria
e reduzir o custo de evolução.

No presente trabalho, a IHM Web emprega a comunicação em tempo real via
WebSocket \cite{fette2011} para a telemetria e a arquitetura REST para comandos e
configuração, uma combinação típica de aplicações modernas de supervisão. A
tabela~\ref{tab:relacionados} resume o posicionamento da solução proposta em
relação às alternativas consideradas.

\begin{table}[htbp]
\centering
\caption{Comparação entre a solução proposta e alternativas para o controle do processo.}
\label{tab:relacionados}
\begin{tabular}{p{3.4cm}p{3.6cm}p{3.6cm}p{3.6cm}}
\toprule
\textbf{Critério} & \textbf{LabVIEW (legado)} & \textbf{CLP dedicado} & \textbf{Solução RPi + Arduino (proposta)} \\
\midrule
Custo de licença/plataforma & Alto & Alto & Baixo \\
Isolamento de tempo real & Parcial & Alto & Alto (Arduino dedicado) \\
Flexibilidade de IHM Web & Baixa & Baixa & Alta \\
Persistência de parâmetros & Limitada & Limitada & Alta (JSON atômico) \\
Watchdog de segurança & Ausente & Variável & Duplo (firmware + backend) \\
Manutenção e evolução & Difícil & Moderada & Facilitada (código aberto) \\
\bottomrule
\end{tabular}\par\vspace{2pt}{\footnotesize Fonte: elaborada pelo autor.}\end{table}

\section{Síntese e lacuna}

A análise dos trabalhos e práticas relacionados indica que, embora existam
soluções comerciais para partes do problema (detecção, controle de temperatura),
não há, no domínio considerado, um sistema aberto e de baixo custo que integre,
de forma coesa e segura, o sequenciamento da preparação de amostras, o controle
PID da rampa do Tubo U e uma IHM Web moderna com persistência de parâmetros. A
lacuna identificada --- e que este trabalho busca preencher --- é exatamente essa
integração, acrescida de mecanismos explícitos de segurança funcional (estado
seguro e watchdog) que o sistema legado não possuía.
